\documentclass[../../main.tex]{subfiles}

\begin{document}

\section{Hilbert-Schmidt 定理}\label{section:ch16-hilbert-schmidt}

非零向量 $u\in H$ 称为是算子 $T\in\mathcal{L}(H)$ 的特征向量, 若存在某个 $\lambda\in\mathbf{R}$ 使得 $T(u)=\lambda u$. 我们称 $\lambda$ 为 $T$ 的与特征向量 $u$ 相应的特征值. 线性代数的中心是以下断言: 若 $H$ 是有限维 Hilbert 空间而 $T\in\mathcal{L}(H)$ 是对称的, 则存在由 $T$ 的特征向量组成的 $H$ 的规范正交基, 即若 $H$ 的维数是 $n$, 存在 $H$ 的规范正交基 $\{\varphi_{1},\ \cdots,\ \varphi_{n}\}$ 以及数 $\{\lambda_{1},\ \cdots,\ \lambda_{n}\}$ 使得对 $1\leqslant k\leqslant n$, $T(\varphi_{k})=\lambda_{k}\varphi_{k}$. 因此,
\begin{align}
\text{对所有 }h\in H,\quad T(h)=\sum_{k=1}^{n}\lambda_{k}\langle h,\ \varphi_{k}\rangle\varphi_{k}.\label{align:ch16-finite-spectral}
\end{align}

当然, 在缺乏对称性的情况下, 有限维空间上的一个有界线性算子, 可能没有任何特征向量. 正如以下例子表明的, 甚至无穷维 Hilbert 空间上的对称算子可以没有任何特征向量.

\begin{example}\label{example:ch16-multiplication-operator}
定义 $T\in\mathcal{L}(L^{2}[a,\ b])$ 为 $[T(f)](x)=xf(x)$, $f\in L^{2}[a,\ b]$. 对 $u,\ v\in L^{2}[a,\ b]$,
\begin{align*}
\langle T(u),v\rangle=\int_{a}^{b}xu(x)v(x)\,dx=\langle u,T(v)\rangle.
\end{align*}
因此 $T$ 是对称的, 且人们容易检验它没有特征向量.
\end{example}

我们将对称算子 $T\in\mathcal{L}(H)$ 与对所有 $h\in H$ 定义为 $Q_{T}(h)=\langle T(h),\ h\rangle$ 的二次型 $Q_{T}:H\to\mathbf{R}$ 相联系. 定义 $T$ 的 Raleigh 商 $R_{T}:H\sim\{0\}\to\mathbf{R}$ 为
\begin{align*}
\text{对所有 }h\in H\sim\{0\},\quad R_{T}(h)=\frac{\langle T(h),h\rangle}{\langle h,h\rangle}
\end{align*}
是有用的.

观察到二次型 $Q_{T}$ 在单位球面 $S=\{h\in H\mid\|h\|=1\}$ 上的最大值点 $h_{*}$ 是 Raleigh 商 $R_{T}$ 在 $H\sim\{0\}$ 上的最大值点.

\begin{lemma}[Hilbert-Schmidt]\label{lemma:ch16-hilbert-schmidt}
令 $H$ 为 Hilbert 空间而 $T\in\mathcal{L}(H)$ 是紧与对称的. 则 $T$ 有特征值 $\lambda$ 使得
\begin{align}
|\lambda|=\|T\|=\sup_{\|h\|=1}|\langle T(h),h\rangle|.\label{align:ch16-hs-lemma}
\end{align}
\end{lemma}

\begin{proof}
若在 $H$ 上 $T=0$, 则 $H$ 中的每个非零向量是 $T$ 的具有特征值 $\lambda=0$ 的特征向量. 因此考虑 $T\neq 0$ 的情形. 命题\cref{proposition:ch16-norm-form}告诉我们,
\begin{align*}
\|T\|=\sup_{\|h\|=1}|\langle T(h),h\rangle|.
\end{align*}
通过用 $-T$ 代替 $T$ 我们可以假定 $\|T\|=\sup_{\|h\|=1}\langle T(h),\ h\rangle$. 记 $\sup_{\|h\|=1}\langle T(h),\ h\rangle$ 为 $\eta$. 令 $S=\{h\in H\mid\|h\|=1\}$ 为 $H$ 的单位球面.

令 $\{h_{k}\}$ 为单位向量序列使得 $\lim_{k\to\infty}\langle T(h_{k}),\ h_{k}\rangle=\eta$. 根据定理\cref{theorem:ch16-weak-seq}, 通过可能地过渡到子序列, 我们可以假定 $\{h_{k}\}$ 弱收敛到 $h_{*}$. 我们有 $\|h_{*}\|\leqslant\liminf\|h_{n}\|=1$, 根据命题\cref{proposition:ch16-compact-weak-strong}, 由于 $T$ 是紧的, $\{T(h_{n})\}$ 强收敛到 $T(h_{*})$. 因此, 根据命题\cref{proposition:ch16-weak-properties},
\begin{align*}
\lim_{k\to\infty}\langle T(h_{k}),h_{k}\rangle=\langle T(h_{*}),h_{*}\rangle.
\end{align*}
因此 $\eta=\langle T(h_{*}),\ h_{*}\rangle$. 现在 $h_{*}\neq 0$, 这是由于 $\eta\neq 0$. 此外, $h_{*}$ 必须是单位向量. 否则 $0<\|h_{*}\|<1$, 在这种情形二次型 $Q_{T}$ 在 $h_{*}/\|h_{*}\|\in S$ 取到大于 $\eta$ 的值, 这与 $\eta$ 是 $Q_{T}$ 在 $S$ 上的上界矛盾. 因此 $h_{*}\in S$ 且对所有 $h\in S$, $Q_{T}(h)\leqslant Q_{T}(h_{*})$. 因此对 $T$ 的 Raleigh 商 $R_{T}$, 我们有
\begin{align*}
\text{对所有 }h\in H\sim\{0\},\quad R_{T}(h)\leqslant R_{T}(h_{*}).
\end{align*}
令 $h_{0}$ 为 $H$ 中的任何向量. 观察到对 $t\in\mathbf{R}$ 定义为 $f(t)=R_{T}(h_{*}+th_{0})$ 的函数 $f:\mathbf{R}\to\mathbf{R}$ 在 $t=0$ 有最大值, 因此 $f^{\prime}(0)=0$. 直接计算给出
\begin{align*}
0=f^{\prime}(0)=\frac{\langle T(h_{0}),h_{*}\rangle+\langle T(h_{*}),h_{0}\rangle}{\|h_{*}\|^{2}}-\langle T(h_{*}),h_{*}\rangle\frac{\langle h_{*},h_{0}\rangle+\langle h_{0},h_{*}\rangle}{\|h_{*}\|^{4}}.
\end{align*}
但 $T$ 是对称的, $h_{*}$ 是单位向量, 且 $\eta=\langle T(h_{*}),\ h_{*}\rangle$, 因此,
\begin{align*}
\langle T(h_{*})-\eta h_{*},h_{0}\rangle=0.
\end{align*}
由于这对所有 $h_{0}\in H$ 成立, $T(h_{*})=\eta h_{*}$.
\end{proof}

\begin{theorem}[Hilbert-Schmidt]\label{theorem:ch16-hilbert-schmidt}
令 $H$ 为 Hilbert 空间. 假定 $K\in\mathcal{L}(H)$ 是紧对称算子且不是有限秩的. 则存在 $[\ker K]^{\perp}$ 的规范正交基 $\{\psi_{k}\}$ 以及非零实数列 $\{\lambda_{k}\}$ 使得 $\lim_{k\to\infty}\lambda_{k}=0$, 且对每个 $k$, $K(\psi_{k})=\lambda_{k}\psi_{k}$. 因此,
\begin{align}
\text{对所有 }h\in H,\quad K(h)=\sum_{k=1}^{\infty}\lambda_{k}\langle h,\ \psi_{k}\rangle\psi_{k}.\label{align:ch16-hs-theorem}
\end{align}
\end{theorem}

\begin{proof}
令 $S$ 为 $H$ 的单位球面. 根据 Hilbert-Schmidt 引理, 我们可以选取向量 $\psi_{1}\in S$ 与 $\mu_{1}\in\mathbf{R}$ 使得
\begin{align*}
K(\psi_{1})=\mu_{1}\psi_{1}\text{ 且 }|\mu_{1}|=\sup_{h\in S}|\langle K(h),h\rangle|.
\end{align*}
由于 $K\neq 0$, 我们从命题\cref{proposition:ch16-norm-form}推出 $\mu_{1}\neq 0$. 定义 $H_{1}=[\operatorname{span}\{\psi_{1}\}]^{\perp}$. 由于 $K(\operatorname{span}\{\psi_{1}\})\subseteq\operatorname{span}\{\psi_{1}\}$, 从命题\cref{proposition:ch16-invariant-orthocomp}得出 $K(H_{1})\subseteq H_{1}$. 因此, 若我们定义 $K_{1}$ 为 $K$ 在 $H_{1}$ 上的限制, 则 $K_{1}\in\mathcal{L}(H_{1})$ 是紧与对称的. 我们再次用 Hilbert-Schmidt 引理选取向量 $\psi_{2}\in S\cap H_{1}$ 与 $\mu_{2}\in\mathbf{R}$ 使得
\begin{align*}
K(\psi_{2})=\mu_{2}\psi_{2}\text{ 且 }|\mu_{2}|=\sup\{\langle K(h),h\rangle\mid h\in S\cap H_{1}\}.
\end{align*}
观察到 $|\mu_{2}|\leqslant|\mu_{1}|$. 此外, 由于 $K$ 没有有限秩, 我们再次用命题\cref{proposition:ch16-norm-form}得出 $\mu_{2}\neq 0$. 我们归纳地选取 $H$ 中的正交的单位向量序列 $\{\psi_{k}\}$ 以及非零实数序列 $\{\mu_{k}\}$, 使得对每个指标 $k$,
\begin{align}
K(\psi_{k})=\mu_{k}\psi_{k}\text{ 且 }|\mu_{k}|=\sup\{|\langle K(h),h\rangle|\mid h\in S\cap[\operatorname{span}\{\psi_{1},\cdots,\psi_{k-1}\}]^{\perp}\}.\label{align:ch16-hs-recursion}
\end{align}
观察到 $\{|\mu_{k}|\}$ 是递减的. 我们宣称 $\{\mu_{k}\}\to 0$. 否则, 由于该序列是递减的, 存在某个 $\varepsilon>0$ 使得对所有 $k$, $|\mu_{k}|\geqslant\varepsilon$. 因此, 对自然数 $m$ 和 $n$, 由于 $\psi_{n}$ 正交于 $\psi_{m}$,
\begin{align*}
\|K(\psi_{n})-K(\psi_{m})\|^{2}=\mu_{n}^{2}\|\psi_{n}\|^{2}+\mu_{m}^{2}\|\psi_{m}\|^{2}\geqslant 2\varepsilon^{2}.
\end{align*}
因此 $\{K(\psi_{k})\}$ 没有强收敛的子序列, 这与算子 $K$ 的紧性矛盾. 因此 $\{\mu_{k}\}\to 0$. 定义 $H_{0}$ 为 $\{\psi_{k}\}_{k=1}^{\infty}$ 的闭线性扩张. 则根据命题\cref{proposition:ch16-complete-iff-basis}, $\{\psi_{k}\}_{k=1}^{\infty}$ 是 $H_{0}$ 的规范正交基. 由于 $K(H_{0})\subseteq H_{0}$, 从命题\cref{proposition:ch16-invariant-orthocomp}得出 $K(H_{0}^{\perp})\subseteq H_{0}^{\perp}$. 但观察到若 $h\in H_{0}^{\perp}$ 是单位向量, 则对每个 $k$, $h\in S\cap[\operatorname{span}\{\psi_{1},\cdots,\psi_{k-1}\}]^{\perp}$, 因此 $|\langle K(h),\ h\rangle|\leqslant|\mu_{k}|$. 由于 $\{\mu_{k}\}\to 0$, $\langle K(h),\ h\rangle=0$, 因此在 $H_{0}^{\perp}$ 上 $Q_{T}=0$, 根据极化等式, $\ker K=H_{0}^{\perp}$. 因此 $H_{0}=[\ker K]^{\perp}$.
\end{proof}

在对称算子 $T\in\mathcal{L}(H)$ 有有限秩的情形, 定义 $H_{0}$ 为 $T$ 的象. 则 $\ker T=H_{0}^{\perp}$. 上面的证明得到了由 $T$ 的特征向量组成的 $H_{0}$ 的规范正交基, 因此重新得到本节开头提到的线性代数的基本结果.

\begin{exercise}\label{exercise:ch16-51}
令 $H$ 为 Hilbert 空间而 $T\in\mathcal{L}(H)$ 为紧与对称的算子. 定义
\begin{align*}
\alpha=\inf_{\|h\|=1}\langle T(h),h\rangle\text{ 与 }\beta=\sup_{\|h\|=1}\langle T(h),h\rangle.
\end{align*}
证明: 若 $\alpha<0$, 则 $\alpha$ 是 $T$ 的特征值; 而若 $\beta>0$, 则 $\beta$ 是 $T$ 的特征值. 给出 $\alpha=0$ 然而 $\alpha$ 不是 $T$ 的特征值的例子, 即 $T$ 是一对一的.
\end{exercise}

\begin{exercise}\label{exercise:ch16-52}
令 $H$ 为 Hilbert 空间而 $K\in\mathcal{L}(H)$ 为紧与对称的. 假定
\begin{align*}
\sup_{\|h\|=1}\langle K(h),h\rangle=\beta>0.
\end{align*}
令 $\{h_{n}\}$ 为单位向量序列使得 $\lim_{n\to\infty}\langle K(h_{n}),\ h_{n}\rangle=\beta$. 证明 $\{h_{n}\}$ 的一个子序列强收敛到 $T$ 的对应于特征值 $\beta$ 的特征向量.
\end{exercise}

\end{document}